% Volume VI: Adversaries, Platforms, and Automated Judgment
% Part XI. Untrustworthy Actors
% Chapter 64: The Asymmetry of Defensive Paranoia
% Spine: proof-spines/ch064-spine.md

\chapter{The Asymmetry of Defensive Paranoia}
\label{ch:ch064}


The same language of suspicion can have opposite political meaning depending on who is authorized to act. The less powerful may need \textbf{defensive paranoia} simply to survive institutions they cannot inspect, predict, or contest. The powerful, by contrast, may invoke the same vocabulary of vigilance to normalize surveillance over people who cannot resist it.

The asymmetry is not merely psychological; it is institutional. Suspicion from below often compensates for opacity and vulnerability, whereas suspicion from above can become an instrument of preemption, monitoring, and control. The chapter's preface therefore insists that paranoia cannot be evaluated in abstraction from power: identical epistemic postures differ morally and politically according to who bears the risk, who commands the apparatus of action, and who can answer back once suspicion is operationalized.
